% sample-kksref.tex
% \kksref（基準ボックス）の実演サンプル。
% コンパイル: latexmk -r .latexmkrc sample-kksref.tex   (LuaLaTeX 専用)
%
% ねらい: 丸囲みのローマ数字小文字 \Rrnum を、\vphantom 手書きなしで
%         一定の見た目に揃える（enumerate のラベルでの典型的な使い方）。
\documentclass{jlreq}
\usepackage[tsumesuji=1]{KKsymbols}
\usepackage{luatexja-otf}
\usepackage{enumitem}

% KKran(\Rrnum) フォールバック（無い環境でも動くように）
\IfFileExists{KKran.sty}{\usepackage{KKran}}{%
  \newcommand{\Rrnum}[1]{\textit{\romannumeral #1}}}

% Hiragino（環境に無ければ適宜差し替え）
\makeatletter
\RequirePackage[no-math]{fontspec}
\RequirePackage[no-math,match,scale=1]{luatexja-fontspec}
\RequirePackage[hiragino-pro,deluxe,expert]{luatexja-preset}
\makeatother

% 丸囲みローマ数字のリスト。ラベルのスコープ内で \kksref を一度設定するだけ。
\newlist{romanmaru}{enumerate}{1}
\setlist[romanmaru,1]{%
  label={\kksref{\Rrnum{8}}\maru{\Rrnum{\arabic*}}},
  leftmargin=3\zw, labelsep=1\zw,
}

\begin{document}

\section*{Before / After}

\noindent\textbf{Before（手書き \texttt{\textbackslash vphantom} が必要だった）:}\par
\maru{\vphantom{\Rrnum{8}}\Rrnum{1}}\,%
\maru{\vphantom{\Rrnum{8}}\Rrnum{4}}\,%
\maru{\vphantom{\Rrnum{8}}\Rrnum{7}}\,%
\maru{\vphantom{\Rrnum{8}}\Rrnum{8}}

\medskip
\noindent\textbf{After（\texttt{\textbackslash kksref} を一度設定するだけ）:}\par
{\kksref{\Rrnum{8}}%
 \maru{\Rrnum{1}}\,\maru{\Rrnum{4}}\,\maru{\Rrnum{7}}\,\maru{\Rrnum{8}}}

\section*{実用例: 丸囲みローマ数字のリスト}

\begin{romanmaru}
  \item 第一の項目。\vphantom は書かなくてよい。
  \item 第二の項目。
  \item 第三の項目。
  \item 第四の項目。すべて同じ大きさ・同じ高さに揃う。
\end{romanmaru}

\end{document}
